% Evaluation methods in reported form; empirical tables are filled from the locked analysis.
\subsection*{Risk sets and paired comparisons}
The unit of evaluation was one disease--outcome--horizon combination. Its risk
set was determined by eligibility and ascertainment of the outcome, not by the
models compared. For each split we report eligible patients, outcome events,
competing deaths and patients whose status was unknown at the horizon. Each
paired comparison used the patients with predictions from both models, so
different contrasts may rest on different risk sets, and their denominators are reported. Cells in which a model could not be fitted or a metric could not be
estimated remain in the coverage inventory with the reason
(\cref{fig:supp-coverage}) and are not treated as zero effects.

Superiority among evaluable cells was kept separate from coverage of
additional cells. A pooled model could predict cells in which no specialist
could be fitted. These cells have absolute performance estimates but no paired
specialist comparison, and were not imputed into pooled averages. Outcome
definitions and the reconstruction of treatment lines are in the appended
data-definition supplement.

\subsection*{Clinical prognostic indices}
Each index was evaluated in the population and for the endpoint closest to
those it was developed for, at the same prediction time. Original and evaluated
endpoints and horizons are listed side by side, and a shorter available
horizon is not presented as a reproduction of a longer development horizon.
All methods in an index comparison used the same patients, those with the index
recorded at diagnosis, and how often it was recorded is reported by disease and
calendar period. Index categories were used as recorded. Because the sequence
and tabular models read the wider history, a gain over an index is a gain of
the whole pipeline and does not separate the model from the information it
receives.

Index categories are not individual probabilities. They were used for ranking
and stratification. Brier scores and calibration of an index would require a
probability mapping fitted on development data, which we did not use.
\placeholder{Index, version, population and availability inventory from the
secure analysis.}

\subsection*{Risk groups, age and operating points}
Continuous model risks were converted to groups fixed at baseline with cutoffs
selected on tuning data, while index groups kept their original definitions.
Kaplan--Meier estimates describe endpoints that include death, and
Aalen--Johansen estimates non-fatal endpoints with death as a competing event.
Group sizes, events and 95\% bootstrap intervals accompany every estimate.
Patients without a recorded index were not included in index comparisons.

In the age figures, patient counts, mean predicted risk and observed risk are
shown as separate rows on the same age bands and groups. The predicted row is
the sequence model's mean probability, also within index categories. The index itself supplies no probability. Estimates are suppressed for cells with fewer
than ten patients, while counts remain visible.

Operating points used thresholds fixed on tuning data, so the specificity
matched on tuning data is not exactly matched in the test period. We report
the achieved sensitivity, specificity, positive and negative predictive values
and patient counts. Confusion counts use patients with known status at the
horizon, and the number excluded by censoring is reported. Censoring-weighted quantities are labelled as such.

For the panel as a whole, every model was also evaluated at six decision
points, each fixed per model, cell and horizon on that model's tuning
predictions. The first flagged as many tuning patients as had the event, so
that every model raised the same number of alerts. The second and third
maximised Youden's index and the MCC. The remaining three were the lowest
thresholds with a tuning specificity of at least 80\%, 90\% and 95\%. The test
period then shows how far the achieved specificity drifts from its tuning
target. For each decision point we report the threshold, the confusion matrix,
sensitivity, specificity, positive and negative predictive values, MCC and the
share of patients flagged, with 95\% intervals from the shared patient
bootstrap and paired differences between the joint model and XGBoost, the
joint model and the specialist, and the routed and joint models. Confusion
counts below five are suppressed in exported tables, as required by the secure
environment. Cells with fewer than five tuning events or non-events have no
decision point.

\subsection*{Censoring, calibration and robustness}
Censoring-weighted and known-status estimates answer different questions and
are reported separately, and the known-status analysis is a prespecified
sensitivity analysis. For the censoring weights we report the censoring model,
the censoring survival at each horizon and the distribution of weights. Long
horizons with little follow-up are flagged where estimates are unstable.

Calibration is reported as calibration-in-the-large, the calibration slope
where estimable, and reliability curves for outcomes chosen by a fixed rule. No
model was recalibrated on test data. Sensitivity analyses by horizon, calendar
year, amount of recorded history and timing of the subtype diagnosis are
reported on their own risk sets. Raw paired effects are reported with the
hierarchical summaries, and the support curve, unsmoothed and binned analyses
are treated as views of the same data rather than independent replications.

\subsection*{Patient representations}
The representation figures show the same test patients at five stages: local
pretraining, the locally pretrained joint model, national pretraining, national
pretraining with local adaptation, and the nationally initialised joint model.
For each stage, the projection was fitted on tuning patients and applied
unchanged to test patients, so coordinates are comparable within a stage but
not between stages. For every stage we record the checkpoint, the layer and
pooling used, the input cutoff and truncation. Clusters and changes in
separation are descriptive. They do not show differences in prediction,
biological mechanisms or the effect of transfer. Only kernel-smoothed
summaries over at least \placeholder{N} patients are shown, following the
disclosure rules of the secure environment.

\subsection*{Numerical tables}
The supplementary tables give, for every cell, the risk sets, absolute
performance, paired effects, index comparisons, calibration, operating points
and sensitivity analyses, together with training resources. Each model was
trained once, so intervals describe sampling of test patients only. Readable
summaries accompany the complete cell-level files.
\placeholder{Populate from the locked analysis.}

\IfFileExists{tables/generated/supplement_tables.tex}{%
  \subsection*{Reviewed numerical summaries}
  \input{tables/generated/supplement_tables}
}{%
  \noindent\placeholder{Reviewed aggregate summary tables have not yet been
  exported. Complete cell-level numerical tables will accompany the figures.}
}
